\BourbakiExercisesSection

\begin{enumerate}
\def\labelenumi{\arabic{enumi}.}
\tightlist
\item
  \begin{enumerate}
  \def\labelenumii{(\alph{enumii})}
  \tightlist
  \item
    If A is a ring with no divisors of zero, show that every element \(\neq0\) of a projective A-module is free.
  \end{enumerate}
\end{enumerate}

\begin{enumerate}
\def\labelenumi{(\alph{enumi})}
\setcounter{enumi}{1}
\item
  Give an example of a projective A-module whose annihilator is non-zero (cf.~I, § 8, no. 11).
\item
  Show that the Z-module Q is not a projective Z-module. *(See Commutative Algebra, II, § 5, Exercise 11 for an example of a finitely generated projective module over an integral domain, which is not free.)*
\end{enumerate}

\begin{enumerate}
\def\labelenumi{\arabic{enumi}.}
\setcounter{enumi}{1}
\item
  Show that every projective A-module is the direct sum of a family of projective submodules each of which admits a countable number of generators (``Kaplansky's Theorem''; apply Exercise 15 of §1, to the case where E is a free module).
\item
  Let \(\mathfrak{c}\) be a cardinal and \(\mathbf{P}\) a projective A-module such that \(\gamma(\mathbf{P}) \leqslant \mathfrak{c}\) (§ 1,
\end{enumerate}

Exercise 6). Show that if L is a free A-module with an infinite basis of cardinal \(\geqslant c\) , P ⊕ L is isomorphic to L. (Observe that there exists a free A-module M, whose basis has cardinal \(\leqslant c\) , isomorphic to the direct sum of P and an A-module Q. Note that M \(^{(N)}\) and P ⊕ M \(^{(N)}\) are isomorphic.)

¶ 4. (a) Let E, E' be two A-modules, N a submodule of E and N' a submodule of E' such that E is projective and E/N and E'/N' are isomorphic. Show that there exists an exact sequence

\[
0 \longrightarrow \mathrm{N} \stackrel {u} {\longrightarrow} \mathrm{E} \oplus \mathrm{N} ^ {\prime} \stackrel {v} {\longrightarrow} \mathrm{E} ^ {\prime} \longrightarrow 0.
\]

(Observe that there exists a homomorphism \(f \colon \mathbf{E} \to \mathbf{E}'\) such that \(f(\mathbf{N}) \subset \mathbf{N}'\) , which gives when passing to the quotients an isomorphism \(\mathbf{E} / \mathbf{N} \to \mathbf{E}' / \mathbf{N}'\) ; define \(u\) and \(v\) in an analogous way to that used for the exact sequence (29) of §1, no. 7.)

\begin{enumerate}
\def\labelenumi{(\alph{enumi})}
\setcounter{enumi}{1}
\item
  Deduce from (a) that if \(\mathbf{E}'\) is projective, \(\mathbf{E} \oplus \mathbf{N}'\) and \(\mathbf{E}' \oplus \mathbf{N}\) are isomorphic.
\item
  Let \(0 \to \mathbf{E}_m \to \mathbf{E}_{m-1} \to \ldots \to \mathbf{E}_1 \to \mathbf{E}_0 \to 0\) be an exact sequence of A-modules such that \(\mathbf{E}_0, \mathbf{E}_1, \ldots, \mathbf{E}_{m-2}\) are projective. Show that the A-modules \(\bigoplus_{h \geqslant 0} \mathbf{E}_{m-2h}\) and \(\bigoplus_{h \geqslant 0} \mathbf{E}_{m-2h+1}\) are isomorphic (with the convention that \(\mathbf{E}_i = \{0\}\) for \(i < 0\) ).
\end{enumerate}

\begin{enumerate}
\def\labelenumi{\arabic{enumi}.}
\setcounter{enumi}{4}
\item
  Let \(u\) be a linear mapping from an A-module \(\mathbf{E}\) to an A-module \(\mathbf{F}\) and \(v = {}^t u\) its transpose. For every submodule \(\mathbf{N}'\) of \(\mathbf{F}^*\) , the submodule of \(\mathbf{E}\) orthogonal to \(v(\mathbf{N}')\) is \(\overline{u}^{-1}(\mathbf{N})\) , where \(\mathbf{N}\) is the submodule of \(\mathbf{F}\) orthogonal to \(\mathbf{N}'\) .
\item
  \begin{enumerate}
  \def\labelenumii{(\alph{enumii})}
  \tightlist
  \item
    If E is an A-module, every injective endomorphism u of E is not a left divisor of zero in the ring \(\operatorname{End}_{\mathbf{A}}(\mathbf{E})\) . Conversely, if for every sub-A-module \(F \neq \{0\}\) of E there exists an endomorphism \(v \neq 0\) of E such that \(v(\mathbf{E}) \subset \mathbf{F}\) , an element which is not a left divisor of zero in \(\operatorname{End}_{\mathbf{A}}(\mathbf{E})\) is an injective endomorphism. The above condition is satisfied if there exists a linear form \(x' \in E^*\) and an \(x \in E\) such that \(\langle x, x' \rangle\) is invertible in A and in particular if E is free.
  \end{enumerate}
\end{enumerate}

\begin{enumerate}
\def\labelenumi{(\alph{enumi})}
\setcounter{enumi}{1}
\item
  Every endomorphism \(\neq 0\) of the \(\mathbf{Z}\) -module \(\mathbf{Q}\) is bijective, although there exists no linear form \(\neq 0\) on \(\mathbf{Q}\) .
\item
  Let \(U_{p}\) be the Z-module defined in §1, Exercise 12(d). Show that the endomorphism \(x \mapsto px\) of \(U_{p}\) is not injective and is not a left divisor of zero in \(\operatorname{End}_{\mathbf{Z}}(\mathbf{U}_{p})\) .
\end{enumerate}

\begin{enumerate}
\def\labelenumi{\arabic{enumi}.}
\setcounter{enumi}{6}
\tightlist
\item
  \begin{enumerate}
  \def\labelenumii{(\alph{enumii})}
  \tightlist
  \item
    If \(u\) is a surjective endomorphism of an A-module \(\mathbf{E}\) , \(u\) is not a right divisor of zero in \(\operatorname{End}_{\mathbf{A}}(\mathbf{E})\) . Conversely, if for every submodule \(\mathbf{F} \neq \mathbf{E}\) of \(\mathbf{E}\) there exists a linear form \(x^{*} \in \mathbf{E}^{*}\) which is zero on \(\mathbf{F}\) and surjective, every element of \(\operatorname{End}_{\mathbf{A}}(\mathbf{E})\) which is not a right divisor of zero is a surjective endomorphism.
  \end{enumerate}
\end{enumerate}

\begin{enumerate}
\def\labelenumi{(\alph{enumi})}
\setcounter{enumi}{1}
\item
  If E is the Z-module defined in Exercise 6(c), show that every endomorphism \(\neq 0\) of E is surjective, although there exists no linear form \(\neq 0\) on E.
\item
  Show that if E is a free Z-module \(\neq0\) , there exist non-surjective endomorphisms of E which are not right divisors of 0 in \(\operatorname{End}_{\mathbf{Z}}(\mathbf{E})\) .
\end{enumerate}

\begin{enumerate}
\def\labelenumi{\arabic{enumi}.}
\setcounter{enumi}{7}
\tightlist
\item
  \begin{enumerate}
  \def\labelenumii{(\alph{enumii})}
  \tightlist
  \item
    Let E be a free A-module, F, G two A-modules and \(u: \mathbf{E} \to \mathbf{G}\) , \(v: \mathbf{F} \to \mathbf{G}\) two linear mappings. Show that if \(u(\mathbf{E}) \subset v(\mathbf{F})\) , there exists a linear mapping \(w: \mathbf{E} \to \mathbf{F}\) such that \(u = v \circ w\) .
  \end{enumerate}
\end{enumerate}

\begin{enumerate}
\def\labelenumi{(\alph{enumi})}
\setcounter{enumi}{1}
\item
  Give an example of two non-zero endomorphisms \(u, v\) of the \(\mathbf{Z}\) -module \(\mathbf{E}\) defined in Exercise 6(c) such that there exists no endomorphism \(w\) of \(\mathbf{E}\) for which \(u = v \circ w\) (although \(u(\mathbf{E}) = v(\mathbf{E}) = \mathbf{E}\) ).
\item
  Give an example of two endomorphisms \(u, v\) of the \(\mathbf{Z}\) -module \(\mathbf{Z}\) such that \(u^{-1}(0) = v^{-1}(0) = \{0\}\) , but such that there exists no endomorphism \(w\) of \(\mathbf{Z}\) for which \(u = w \circ v\) (cf.~§ 7, Exercise 14).
\end{enumerate}

\begin{enumerate}
\def\labelenumi{\arabic{enumi}.}
\setcounter{enumi}{8}
\tightlist
\item
  Given an A-module E, for every submodule M of E (resp. every submodule M' of E*), let M° (resp. M'°) denote the orthogonal of M in E* (resp. the orthogonal of M' in E). Consider the following four properties:
\end{enumerate}

\begin{enumerate}
\def\labelenumi{(\Alph{enumi})}
\item
  The canonical homomorphism \(c_{\mathbf{E}}: \mathbf{E} \to \mathbf{E}^{**}\) is bijective.
\item
  For every submodule M of E, the canonical homomorphism \(E^{*}/M^{\circ} \rightarrow M^{*}\) is bijective.
\item
  For every submodule M of E, \(M^{\circ\circ} = M\) .
\item
  For every ordered pair of submodules M, N of E,
\end{enumerate}

\[
(\mathbf {M} \cap \mathbf {N}) ^ {\circ} = \mathbf {M} ^ {\circ} + \mathbf {N} ^ {\circ}.
\]

\begin{enumerate}
\def\labelenumi{(\alph{enumi})}
\item
  Show that condition (B) implies (D) (prove that for \(x^{*} \in (\mathbf{M} \cap \mathbf{N})^{\circ}\) , there exists \(y^{*} \in \mathbf{E}^{*}\) such that \(\langle x + y, y^{*} \rangle = \langle x, x^{*} \rangle\) for \(x \in \mathbf{M}\) and \(y \in \mathbf{N}\) ).
\item
  Let A be a ring with no divisors of zero, but having no field of left fractions \(*\) (for example the tensor algebra of a vector space of dimension \textgreater1, cf.~III, § 5, Exercise 5) \(^{*}\) . If we take \(E = A_{s}\) , condition (A) is satisfied, but none of conditions (B), (C), (D) (for an example where (B), (C), (D) hold, but not (A), see § 7, no. 5, Theorem 6).
\item
  Take A to be a product ring \(\mathbf{K}^{\mathrm{I}}\) , where \(\mathbf{K}\) is a commutative field and I an infinite set. Show that the A-module \(\mathbf{E} = \mathbf{A}\) satisfies conditions (A) and (B), but not (C) (observe that the annihilators of the ideals of \(\mathbf{A}\) are all of the form \(\mathbf{K}^{\mathrm{J}}\) , where \(\mathbf{J} \subset \mathbf{I}\) ).
\item
  Suppose that for two submodules \(\mathbf{M},\mathbf{N}\) of \(\mathbf{E}\) such that \(\mathbf{N} \subset \mathbf{M}\) and \(\mathbf{N} \neq \mathbf{M}\) , the dual \((\mathbf{M} / \mathbf{N})^*\) is non-zero; then show that condition (B) implies (C).
\item
  The kernel of \(c_{\mathbf{E}}\) is the orthogonal \((\mathbf{E}^{*})^{\circ}\) of \(\mathbf{E}^{*}\) in \(\mathbf{E}\) . Give an example where neither \(\mathbf{E}^{*}\) nor \((\mathbf{E}^{*})^{\circ}\) is 0 (consider a module containing an element whose annihilator contains an element which is not a divisor of 0).
\item
  Let \(\mathbf{M}\) be a direct factor of \(\mathbf{E}\) ; show that \(\mathbf{M}^{\circ \circ} = \mathbf{M} + (\mathbf{E}^{*})^{\circ}\) .
\end{enumerate}

\begin{enumerate}
\def\labelenumi{\arabic{enumi}.}
\setcounter{enumi}{9}
\tightlist
\item
  Give an example of an A-linear mapping \(u: \mathbf{E} \to \mathbf{F}\) such that \(^t u\) is bijective but \(u\) is neither injective nor surjective (cf.~Exercise 6(b) and (c)).
\end{enumerate}

Deduce an example where \(y_{0} \in F\) is orthogonal to the kernel of \(^{t}u\) but where the equation \(u(x) = y_{0}\) has no solution.

¶ 11. Let A be a ring and I an A-module. I is called injective if, for every exact sequence \(\mathrm{E}' \to \mathrm{E} \to \mathrm{E}''\) of A-linear mappings, the sequence

\[
\operatorname{Hom} \left(\mathrm{E} ^ {\prime \prime}, \mathrm{I}\right)\rightarrow \operatorname{Hom} (\mathrm{E}, \mathrm{I}) \rightarrow \operatorname{Hom} \left(\mathrm{E} ^ {\prime}, \mathrm{I}\right)
\]

is exact.

\begin{enumerate}
\def\labelenumi{(\alph{enumi})}
\tightlist
\item
  Show that the following properties are equivalent:
\end{enumerate}

\((\alpha)\) I is injective.

(β) For every exact sequence \(0 \rightarrow E' \rightarrow E \rightarrow E'' \rightarrow 0\) of A-linear mappings, the sequence

\[
0 \rightarrow \operatorname{Hom} (\mathrm{E} ^ {\prime \prime}, \mathrm{I}) \rightarrow \operatorname{Hom} (\mathrm{E}, \mathrm{I}) \rightarrow \operatorname{Hom} (\mathrm{E} ^ {\prime}, \mathrm{I}) \rightarrow 0
\]

is exact.

(γ) For every A-module M and every submodule N of M, every A-linear mapping of N into I can be extended to an A-linear mapping of M into I.

(δ) For every left ideal \(\mathfrak{a}\) of A and every A-linear mapping \(f: \mathfrak{a} \to \mathbf{I}\) , there exists an element \(b \in \mathbf{I}\) such that \(f(a) = ab\) for all \(a \in \mathfrak{a}\) .

(ζ) I is a direct factor of every A-module containing it.

(θ) For every A-module E which is the sum of I and a monogenous submodule, I is a direct factor of E.

(To prove that (δ) implies (γ), show that property (δ) implies that if E is an A-module and F a submodule of E such that E/F is monogenous, then every linear mapping of F into I can be extended to a linear mapping of E into I. Then use Zorn's Lemma. To prove that (θ) implies (δ), consider the A-module \(A_{s} \times I = M\) , the submodule N of M consisting of the elements \((a, -f(a))\) for \(a \in \mathfrak{a}\) , and apply (θ) to the quotient module M/N.)

\begin{enumerate}
\def\labelenumi{(\alph{enumi})}
\setcounter{enumi}{1}
\tightlist
\item
  For a Z-module E to be injective, it is necessary and sufficient that for all \(x \in E\) and every integer \(n \neq 0\) , there exists \(y \in E\) such that ny = x. In particular, the Z-modules Q and Q/Z are injective.
\end{enumerate}

\begin{enumerate}
\def\labelenumi{\arabic{enumi}.}
\setcounter{enumi}{11}
\tightlist
\item
  \begin{enumerate}
  \def\labelenumii{(\alph{enumii})}
  \tightlist
  \item
    For a product \(\prod_{\lambda \in L} E_{\lambda}\) of A-modules to be injective (Exercise 11), it is necessary and sufficient that each of the \(E_{\lambda}\) be injective.
  \end{enumerate}
\end{enumerate}

\begin{enumerate}
\def\labelenumi{(\alph{enumi})}
\setcounter{enumi}{1}
\tightlist
\item
  Let K be a commutative field, L an infinite set and A the product ring \(K^{L}\) . Show that A is an injective A-module (use (a)) but that the ideal \(\mathfrak{a}\) of A the direct sum of the factors of \(K^{L}\) (canonically identified with ideals of A) is not an injective A-module.
\end{enumerate}

¶ 13. (a) Let A be a ring and I an injective A-module (Exercise 11) such that for every non-zero monogenous A-module E, there exists a non-zero A-homomorphism of E into I. Let M be an A-module, N a submodule of M and \(\tilde{\mathbf{N}}\) the set of \(u \in \operatorname{Hom}_{\mathbf{A}}(\mathbf{M}, \mathbf{I})\) such that \(u(x) = 0\) in N. Show that for all \(y \notin N\) , there exists a \(u \in \tilde{\mathbf{N}}\) such that \(u(y) \neq 0\) .

\begin{enumerate}
\def\labelenumi{(\alph{enumi})}
\setcounter{enumi}{1}
\item
  Let A be a ring. For every (A, A)-bimodule T and every left (resp. right) A-module M, \(\operatorname{Hom}_{\mathbf{A}}(\mathbf{M}, \mathbf{T})\) has canonically a right (resp. left) A-module structure deriving from that on T, and there is a canonical A-homomorphism \(c_{M,T}: M \to \operatorname{Hom}_{\mathbf{A}}(\operatorname{Hom}_{\mathbf{A}}(\mathbf{M}, \mathbf{T}), \mathbf{T})\) such that \(c_{M,T}(x)\) is the A-homomorphism \(u \mapsto u(x)\) . Show that if I is an (A, A)-bimodule which, as a left A-module, satisfies the conditions of (a), then the canonical homomorphism \(c_{M,I}\) is injective for every left A-module M.
\item
  Suppose that I is an (A, A)-bimodule which is injective as a left A-module and a right A-module and that for every (left or right) monogenous A-module \(E \neq 0\) , there exists a non-zero A-homomorphism of E into I. Show that under these conditions, if P is a left (resp. right) projective A-module, \(\operatorname{Hom}_{\mathbf{A}}(\mathbf{P}, \mathbf{I})\) is a right (resp. left) injective A-module. (Suppose that P is a left projective A-module. Observe that for every submodule \(N'\) of a right A-module N, \(\operatorname{Hom}_{\mathbf{A}}(\mathbf{N}', \mathbf{I})\) is isomorphic to a quotient module of \(\operatorname{Hom}_{\mathbf{A}}(\mathbf{N}, \mathbf{I})\) and apply (b) to the right module N and the left module P.)
\item
  Let C be a commutative ring, A a C-algebra and I an injective C-module such that for every monogenous C-module \(E \neq \{0\}\) there exists a non-zero C-homomorphism of E into I. For every A-module M, \(\operatorname{Hom}_{\mathbf{C}}(\mathbf{M}, \mathbf{I})\) is an A-module; show that for every projective A-module P, \(\operatorname{Hom}_{\mathbf{C}}(\mathbf{P}, \mathbf{I})\) is an injective A-module (same argument as in (c)).
\end{enumerate}

\begin{enumerate}
\def\labelenumi{\arabic{enumi}.}
\setcounter{enumi}{13}
\tightlist
\item
  Let A be a ring. Show that for every A-module M there exists an injective A-module I such that M is isomorphic to a submodule of I. (Apply Exercise 13(d) with C = Z, using Exercise 11(b) and the fact that every A-module is the quotient of a free A-module.)
\end{enumerate}

¶ 15. An A-module homomorphism \(u: \mathbf{E} \to \mathbf{F}\) is called essential if it is injective and if, for every submodule \(\mathbf{P} \neq \{0\}\) of \(\mathbf{F}\) , \(u^{-1}(\mathbf{P}) \neq \{0\}\) ; it suffices that this condition holds for every monogenous submodule \(\mathbf{P} \neq \{0\}\) of \(\mathbf{F}\) . If \(\mathbf{E}\) is a submodule of \(\mathbf{F}\) , \(\mathbf{F}\) is called an essential extension of \(\mathbf{E}\) if the canonical injection \(\mathbf{E} \to \mathbf{F}\) is essential.

\begin{enumerate}
\def\labelenumi{(\alph{enumi})}
\item
  Let \(u: \mathbf{E} \to \mathbf{F}\) , \(v: \mathbf{F} \to \mathbf{G}\) be two A-homomorphisms of A-modules. If \(u\) and \(v\) are essential, so is \(v \circ u\) . Conversely, if \(v \circ u\) is essential and \(v\) is injective, then \(u\) and \(v\) are essential.
\item
  Let E be a submodule of an A-module F and let \((\mathbf{F}_{\lambda})_{\lambda \in L}\) be a family of submodules of F containing E and whose union is F. Show that if each of the \(F_{\lambda}\) is an essential extension of E, so is F.
\item
  Let \((u_{\lambda})_{\lambda \in L}\) be a family of essential homomorphisms \(u_{\lambda} \colon \mathbf{E}_{\lambda} \to \mathbf{F}_{\lambda}\) . Show that the homomorphism \(\bigoplus_{\lambda \in L} u_{\lambda} \colon \bigoplus_{\lambda \in L} \mathbf{E}_{\lambda} \to \bigoplus_{\lambda \in L} \mathbf{F}_{\lambda}\) is essential. (Prove it first when \(L\) has two elements, then use (b).)
\item
  The \(\mathbf{Z}\) -module \(\mathbf{Q}\) is an essential extension of \(\mathbf{Z}\) but the product \(\mathbf{Q}^{\mathrm{N}}\) is not an essential extension of \(\mathbf{Z}^{\mathrm{N}}\) .
\item
  Let E be a submodule of a module F. Show that there exists a submodule
\end{enumerate}

Q of F such that \(Q \cap E = \{0\}\) and that the restriction to E of the canonical homomorphism \(F \to F/Q\) is essential.

\begin{enumerate}
\def\labelenumi{\arabic{enumi}.}
\setcounter{enumi}{15}
\tightlist
\item
  A submodule E of an A-module F is called irreducible with respect to F (or in F) if \(E \neq F\) and E is not the intersection of two submodules of F distinct from E.
\end{enumerate}

\begin{enumerate}
\def\labelenumi{(\alph{enumi})}
\item
  For an A-module F to be an essential extension of each of its submodules \(\neq \{0\}\) it is necessary and sufficient that \(\{0\}\) be irreducible with respect to F.
\item
  Let \((\mathbf{E}_{\lambda})_{\lambda \in L}\) be a family of submodules of an A-module \(\mathbf{F}\) . \((\mathbf{E}_{\lambda})_{\lambda \in L}\) is called a reduced irreducible decomposition of \(\{0\}\) in \(\mathbf{F}\) if each of the \(\mathbf{F}_{\lambda}\) is irreducible in \(\mathbf{F}\) , if the intersection of the \(\mathbf{E}_{\lambda}\) reduces to 0 and if none of the \(\mathbf{E}_{\lambda}\) contains the intersection of the \(\mathbf{E}_{\mu}\) for \(\mu \neq \lambda\) . When this is so, show that the canonical homomorphism of \(\mathbf{F}\) into \(\bigoplus_{\lambda \in L} (\mathbf{F} / \mathbf{E}_{\lambda})\) is essential (if \(\mathbf{E}_{\lambda}' = \bigcap_{\mu \neq \lambda} \mathbf{E}_{\mu}\) , observe that the image \(\mathbf{F}_{\lambda}'\) of \(\mathbf{E}_{\lambda}'\) in \(\mathbf{F} / \mathbf{E}_{\lambda}\) is non-zero and that the image of \(\mathbf{F}\) in \(\bigoplus_{\lambda \in L} (\mathbf{F} / \mathbf{E}_{\lambda})\) contains \(\bigoplus_{\lambda \in L} \mathbf{F}_{\lambda}'\) ; then apply Exercise 15(c)). Conversely, if \((\mathbf{E}_{\lambda})_{\lambda \in L}\) is a family of submodules of \(\mathbf{F}\) irreducible in \(\mathbf{F}\) and the canonical homomorphism \(\mathbf{F} \to \bigoplus_{\lambda \in L} (\mathbf{F} / \mathbf{E}_{\lambda})\) is essential, the family \((\mathbf{E}_{\lambda})\) is a reduced irreducible decomposition of \(\{0\}\) in \(\mathbf{F}\) .
\end{enumerate}

¶ 17. (a) Let M, N, M', N' be four submodules of a module E such that M ∩ N = M' ∩ N' = P. Show that

\[
\mathbf {M} = (\mathbf {M} + (\mathbf {N} \cap \mathbf {M} ^ {\prime})) \cap (\mathbf {M} + (\mathbf {N} \cap \mathbf {N} ^ {\prime})).
\]

Deduce that if \(\mathbf{M}\) is irreducible in \(\mathbf{E}\) , then necessarily \(\mathbf{P} = \mathbf{M}' \cap \mathbf{N}\) or \(\mathbf{P} = \mathbf{N}' \cap \mathbf{N}\) .

\begin{enumerate}
\def\labelenumi{(\alph{enumi})}
\setcounter{enumi}{1}
\tightlist
\item
  Let \((\mathrm{N}_{i})_{1\leqslant i\leqslant n}\) be a finite family of submodules of E irreducible in E and M its intersection. Show that if \((\mathrm{P}_{j})_{1\leqslant j\leqslant m}\) is a finite family of submodules of E such that \(M=\bigcap_{j}P_{j}\) , then for every index i such that \(1\leqslant i\leqslant n\) there exists an index \(\phi(i)\) such that \(1\leqslant\phi(i)\leqslant m\) and such that, if we write \(N_{i}^{\prime}=\bigcap_{k\neq i}N_{k}\) , then \(M=P_{\phi(i)}\cap N_{i}^{\prime}\) (use (a)). Deduce that if none of the \(P_{j}\) contains the intersection of the \(P_{k}\) of index \(\neq j\) , then \(m\leqslant n\) (successively replace the \(N_{i}\) by suitable \(P_{j}\) , using the above result). Conclude that two reduced irreducible decompositions of M in E, one of which is finite, necessarily have the same number of terms (cf.~Exercise 23).
\end{enumerate}

¶ 18. (a) For an A-module I to be injective, show that it is necessary and sufficient that I admit no essential extension distinct from itself (show that this condition implies condition (δ) of Exercise 11(a), arguing as for the proof that (θ) implies (δ) in that exercise).

\begin{enumerate}
\def\labelenumi{(\alph{enumi})}
\setcounter{enumi}{1}
\item
  Let I be an injective A-module. For a submodule E of I to be injective, it is necessary and sufficient that E admit no essential extension distinct from itself and contained in I (using Exercise 15(e)), show that this condition implies that E is a direct factor of I).
\item
  Let I be an injective A-module, E a submodule of I and \(h: \mathbf{E} \to \mathbf{F}\) an essential homomorphism. Show that there exists an injective homomorphism \(j: \mathbf{F} \to \mathbf{I}\) such that \(j \circ h\) is the canonical injection of E into I.
\item
  Let I be an A-module and E a submodule of I. Show that the following properties are equivalent:
\end{enumerate}

(α) I is injective and is an essential extension of E.

(β) I is injective and for every injective homomorphism \(j\) of \(E\) into an injective A-module \(J\) , there exists an injective homomorphism of \(I\) into \(J\) extending \(j\) .

(γ) I is injective and is the smallest injective submodule of I containing E.

(δ) I is an essential extension of E and for every essential homomorphism \(h: E \to F\) , there exists an injective homomorphism \(j: F \to I\) such that \(j \circ h\) is the canonical injection of E into I.

(Use (c) and (a).) When these equivalent conditions hold, I is called an injective envelope of E; I is then also an injective envelope of every submodule of I containing E.

\begin{enumerate}
\def\labelenumi{(\alph{enumi})}
\setcounter{enumi}{4}
\item
  Let I be an injective A-module and E a submodule of I. Show that every maximal element of the set of essential extensions of E contained in I is an injective envelope of E (use (b)). In particular, every A-module admits an injective envelope (cf.~Exercise 14).
\item
  If I, I' are two injective envelopes of the same A-module E, show that there exists an isomorphism of I onto I' leaving invariant the elements of E.
\end{enumerate}

¶ 19. (a) Let E, F be two A-modules and M an injective submodule of E ⊕ F. Let I be an injective envelope of M ∩ E in M and J a supplementary submodule of I in M. Show that the restriction to I (resp. J) of the canonical projection of E ⊕ F onto E (resp. F) is injective (compose with the projection of I into E an extension E → I of the injection M ∩ E → I).

\begin{enumerate}
\def\labelenumi{(\alph{enumi})}
\setcounter{enumi}{1}
\tightlist
\item
  Let E be an A-module and M a submodule of E, maximal in the set of injective submodules of E; M admits in E a supplementary submodule N which contains no injective submodule \(\neq\{0\}\) . Show that for every injective submodule P of E, the image of P under the projection of E onto M (under the decomposition of E as a direct sum \(M \oplus N\) ) is an injective envelope of \(P \cap M\) in M (use (a)). If \(M'\) is another submodule of E which is maximal in the set of injective submodules of E, show that there exists an automorphism of E transforming M into \(M'\) and leaving invariant the elements of N.
\end{enumerate}

¶ 20. (a) Let A be a ring admitting a field of left fractions K (I, § 9, Exercise 15). Show that K, considered as a left A-module, is an injective envelope of A \(_{s}\) . (Apply criterion (α) of Exercise 18(d) and criterion (δ) of Exercise 11(a) noting that if f: a → K is an A-linear mapping of a left ideal a of A into K, the element \(x^{-1}f(x)\) for \(x \in \mathfrak{a}, x \neq 0\) (the inverse being taken in K) does not depend on \(x \in \mathfrak{a}\) .)

\begin{enumerate}
\def\labelenumi{(\alph{enumi})}
\setcounter{enumi}{1}
\tightlist
\item
  Let \(U_{p}\) be the sub-Z-module of Q/Z consisting of the canonical images of the rational numbers of the form \(k/p^{n}\) , where p is a given prime number, \(k \in Z\) and \(n \in N\) (II, §1, Exercise 12(d)). Show that \(U_{p}\) is an injective envelope of each of the Z-modules \(p^{-n}Z/Z\) isomorphic to \(Z/p^{n}Z\) . Show that there exist automorphisms of \(U_{p}\) distinct from the identity automorphism leaving invariant the elements of a submodule \(p^{-n}Z/Z\) .
\end{enumerate}

¶ 21. (a) Let \((\mathbf{E}_j)_{j \in J}\) be a finite family of A-modules and for all \(j \in J\) let \(I_j\) be an injective envelope of \(E_j\) ; show that \(\bigoplus_{j \in J} I_j\) is an injective envelope of \(\bigoplus_{j \in J} E_j\) .

\begin{enumerate}
\def\labelenumi{(\alph{enumi})}
\setcounter{enumi}{1}
\tightlist
\item
  An A-module M is called indecomposable if it is distinct from \(\{0\}\) and it is not the direct sum of two submodules distinct from \(\{0\}\) . Show that if I is an injective A-module, the following properties are equivalent:
\end{enumerate}

\((\alpha)\) \(\{0\}\) is an irreducible submodule in I.

(β) I is indecomposable.

(γ) I is an injective envelope of each of its submodules \(\neq0\) . (Use (a) and Exercise 18 (e).)

(δ) I is isomorphic to the injective envelope \(\mathrm{I}(\mathrm{A}/\mathrm{q})\) , where q is an irreducible left ideal of A.

Moreover, when this is so, for all \(x \neq 0\) in I, \(\operatorname{Ann}(x)\) is an irreducible left ideal of A and I is isomorphic to \(\operatorname{I}(\mathbf{A} / \operatorname{Ann}(x))\) .

Deduce that for the injective envelope of an A-module E to be indecomposable, it is necessary and sufficient that \(\{0\}\) be an irreducible submodule of E.

\begin{enumerate}
\def\labelenumi{(\alph{enumi})}
\setcounter{enumi}{2}
\item
  Give an example of an indecomposable Z-module E such that its injective envelope I(E) is decomposable (cf.~VII, § 3, Exercise 5).
\item
  Let E be an A-module, I an injective envelope of E and \(I = \bigoplus_{\lambda \in L} I_{\lambda}\) a decomposition of I as a direct sum of a finite family of indecomposable injective submodules; for all \(\lambda \in L\) let \(J_{\lambda} = \bigoplus_{\mu \neq \lambda} I_{\mu}\) and let \(N_{\lambda} = J_{\lambda} \cap E\) . Show that \((\mathrm{N}_{\lambda})_{\lambda \in L}\) is a reduced irreducible decomposition of \(\{0\}\) in E (Exercise 16 (b)) and that \(I_{\lambda}\) is an injective envelope of \(E/N_{\lambda}\) . Conversely, every reduced irreducible decomposition of \(\{0\}\) in E can be obtained by the above procedure uniquely up to isomorphism (use Exercise 16 (b)).
\end{enumerate}

\begin{enumerate}
\def\labelenumi{\arabic{enumi}.}
\setcounter{enumi}{21}
\tightlist
\item
  Let I be an injective A-module. If I is indecomposable, every injective endomorphism of I is an automorphism of I. Deduce that, for I to be indecomposable, it is necessary and sufficient that the non-invertible elements of the ring End(I) form a two-sided ideal of that ring.
\end{enumerate}

¶ 23. Let M be an A-module, the direct sum of a family (finite or otherwise) \((\mathbf{M}_{\lambda})_{\lambda \in L}\) of submodules such that for all \(\lambda \in L\) the non-invertible elements of \(\operatorname{End}(\mathbf{M}_{\lambda})\) form a two-sided ideal of this ring (which implies that \(\mathbf{M}_{\lambda}\) is indecomposable).

\begin{enumerate}
\def\labelenumi{(\alph{enumi})}
\item
  Let \(f, g\) be two endomorphisms of \(\mathbf{M}\) such that \(f + g = 1_{\mathbf{M}}\) . Show that for every finite sequence \((\lambda_k)_{1 \leqslant k \leqslant s}\) of distinct indices in \(\mathbf{L}\) there exists a family of submodules \(\mathbf{N}_k\) ( \(1 \leqslant k \leqslant s\) ) of \(\mathbf{M}\) such that for every \(k\) one of the endomorphisms \(f, g\) restricted to \(\mathbf{M}_{\lambda_{k}}\) is an isomorphism of this submodule onto \(\mathbf{N}_k\) and that \(\mathbf{M}\) is the direct sum of the \(\mathbf{N}_k\) ( \(1 \leqslant k \leqslant s\) ) and the \(\mathbf{M}_{\lambda}\) for \(\lambda\) distinct from \(\lambda_k\) . (Reduce it to the case \(s = 1\) ; if \(\pi_{\lambda}\) is the canonical projection of \(\mathbf{M}\) onto \(\mathbf{M}_{\lambda}\) corresponding to the decomposition as a direct sum of \(\mathbf{M}_{\lambda}\) and \(\bigoplus_{\mu \neq \lambda} \mathbf{M}_{\mu}\) , note that either \(\pi_{\lambda} \circ f\) or \(\pi_{\lambda} \circ g\) , restricted to \(\mathbf{M}_{\lambda}\) , is an automorphism of \(\mathbf{M}_{\lambda}\) .)
\item
  Let \(f\) be an idempotent endomorphism of \(\mathbf{M}\) . Show that there exists at least one index \(\lambda \in \mathbf{L}\) such that the restriction of \(f\) to \(\mathbf{M}_{\lambda}\) is an isomorphism of \(\mathbf{M}_{\lambda}\) onto \(f(\mathbf{M}_{\lambda})\) ; moreover \(f(\mathbf{M}_{\lambda})\) is a direct factor of \(\mathbf{M}\) . (If \(g = 1_{\mathbf{M}} - f\) , note that there exists a finite sequence of indices \((\lambda_k)_{1 \leqslant k \leqslant s}\) such that the intersection of \(\operatorname{Ker}(g)\) and the direct sum of the \(\mathbf{M}_{\lambda_k}\) is \(\neq \{0\}\) and use (a).)
\item
  Deduce from (b) that every indecomposable direct factor of M is isomorphic to one of the \(M_{\lambda}\) .
\item
  Let \((\mathrm{N}_{\kappa})_{\kappa\in K}\) be a family of indecomposable submodules of M of which M is the direct sum; every \(N_{\kappa}\) is therefore isomorphic to an \(M_{\lambda}\) and vice versa, by virtue of (c). Let T be the set of classes of indecomposable submodules of M (under the relation of isomorphism) such that an \(M_{\lambda}\) (or an \(N_{\kappa}\) ) belongs to one of these classes. For all \(t\in T\) , let \(\mathrm{R}(t)\) (resp. \(\mathrm{S}(t)\) ) be the set of \(\lambda\in L\) (resp. \(\kappa\in K\) ) such that \(M_{\lambda}\in t\) (resp. \(N_{\kappa}\in t\) ). Show that, for all \(t\in T\) , \(\mathrm{Card}(\mathrm{R}(t))=\mathrm{Card}(\mathrm{S}(t))\) . (In the notation of (a), let \(\mathrm{J}(\kappa)\) be the set of \(\lambda\in L\) such that the restriction of \(\pi_{\lambda}\) to \(N_{\kappa}\) is an isomorphism of \(N_{\kappa}\) onto \(M_{\lambda}\) ; show that \(\mathrm{J}(\kappa)\) is finite and that the \(\mathrm{J}(\kappa)\) cover \(\mathrm{R}(t)\) when \(\kappa\) runs through \(\mathrm{S}(t)\) . Deduce that \(\mathrm{Card}(\mathrm{S}(t))\leqslant\mathrm{Card}(\mathrm{R}(t))\) when \(\mathrm{R}(t)\) is infinite. When \(\mathrm{R}(t)\) is finite, show with the aid of (a) that M is the direct sum of the \(M_{\lambda}\) for \(\lambda\notin R(t)\) and a subfamily of \((\mathrm{N}_{\kappa})_{\kappa\in\mathrm{S}(t)}\) of cardinal equal to that of \(\mathrm{R}(t)\) .
\item
  Deduce from (d) and Exercises 22 and 21 (d) that if \((\mathbf{E}_{\lambda})_{\lambda \in L}\) and \((\mathbf{E}_{\kappa}^{\prime})_{\kappa \in K}\) are two reduced irreducible decompositions of \(\{0\}\) in a module F, L and K are equipotent.
\end{enumerate}

\begin{enumerate}
\def\labelenumi{\arabic{enumi}.}
\setcounter{enumi}{23}
\item
  Let M be an A-module, I its injective envelope (Exercise 18), \(\mathfrak{a}\) a two-sided ideal of A and Q the submodule of I consisting of the \(z \in I\) such that \(\mathfrak{a}z = \{0\}\) ; Q has a natural \((\mathrm{A}/\mathfrak{a})\) -module structure. Let N be the submodule of M consisting of the \(x \in M\) such that \(\mathfrak{a}x = \{0\}\) ; if N is considered as an \((\mathrm{A}/\mathfrak{a})\) -module, show that Q is isomorphic to an injective envelope of N.
\item
  \begin{enumerate}
  \def\labelenumii{(\alph{enumii})}
  \tightlist
  \item
    For an A-module \(\mathbf{Q}\) to be injective, it is sufficient that, for every projective A-module P and every submodule \(P'\) of P, every linear mapping of \(P'\) into Q can be extended to a linear mapping of P into Q (use the fact that every A-module is a quotient of a projective A-module).
  \end{enumerate}
\end{enumerate}

\begin{enumerate}
\def\labelenumi{(\alph{enumi})}
\setcounter{enumi}{1}
\tightlist
\item
  For an A-module P to be projective, it is sufficient that, for every injective A-module Q and every quotient module \(Q''\) of Q, every homomorphism of P into \(Q''\) be of the form \(\phi \circ u\) , where \(\phi: Q \to Q''\) is the canonical homomorphism and u is a homomorphism of P into Q (use the fact that every A-module is a submodule of an injective A-module).
\end{enumerate}

\begin{enumerate}
\def\labelenumi{\arabic{enumi}.}
\setcounter{enumi}{25}
\tightlist
\item
  For every Z-module G and every integer n \textgreater{} 0, let \(_{n}G\) denote the kernel of the endomorphism \(x \mapsto nx\) of G. If \(0 \to G' \to G \to G'' \to 0\) is an exact sequence, define a canonical homomorphism \(d:_{n}G'' \to G'/nG'\) such that the sequence
\end{enumerate}

\[
0 \longrightarrow_ {n} G ^ {\prime} \longrightarrow_ {n} G ^ {\prime \prime} \stackrel {d} {\longrightarrow} G ^ {\prime} / n G ^ {\prime} \longrightarrow G / n G \longrightarrow G ^ {\prime \prime} / n G ^ {\prime \prime} \longrightarrow 0
\]

is exact. If \(_n\mathrm{G}\) and \(\mathrm{G} / n\mathrm{G}\) are finite, we write

\[
h _ {n} (\mathrm{G}) = \operatorname{Card} (\mathrm{G} / n \mathrm{G}) - \operatorname{Card} (_ {n} \mathrm{G});
\]

then show that \(_{n}G'\) , \(_{n}G''\) , \(G'/nG'\) , \(G''/nG''\) are finite and that

\[
h _ {n} (\mathrm{G}) = h _ {n} \left(\mathrm{G} ^ {\prime}\right) + h _ {n} \left(\mathrm{G} ^ {\prime \prime}\right).
\]

